\section{Related Work}
\label{sec:related}
\subsection{Industrial Agents and Plant-Facing Action}
Industrial agents supervise and reconfigure service-oriented cyber--physical systems
\cite{biskupovic2025supervision}. Recent LLM work extends this direction to
PLC generation and verification \cite{liu2026agents4plc}, instruction-guided
control \cite{wu2025instructmpc}, digital-twin fault handling
\cite{gill2025llmdt}, and natural-language machine control over OPC UA
\cite{hofmann2025nlc}. These studies motivate a reusable boundary between an
agent's proposal and its industrial execution.

The closest plant-facing systems emphasize complementary functions
(Table~\ref{tab:nearest}): operator-reviewed control-sequence generation with
OPC UA execution \cite{huang2026made}, module-scoped skill dispatch
\cite{koehle2026llmdriven}, and LLM scheduling in a simulated production
line \cite{waseem2026pretrained}. Agent-allocation systems emphasize resource
capability exploration and shop-floor coordination
\cite{lim2025llmenabled,zhao2026shopfloor}. Our comparison follows the
authority and execution evidence connecting these strategies to production
objects. The framework's approval and capability boundary implements the
safety-aware human-in-the-loop decision-support perspective discussed in
\cite{jiang2026agentic}.

K{\"o}hle et al.\ also restrict tool availability by component state and
cross-verify states to expose silent faults in a simulated six-module plant
\cite{koehle2026llmdriven}. These mechanisms overlap with our admission and
confirmation goals. The present design organizes that relationship around
production-scoped capabilities, source-specific limits, and separate
receipt/verdict semantics. The integration links three contracts:
production-scoped authorization, graded command evidence, and separate
command, observation, and recipe records. Table~\ref{tab:nearest} positions
this combined contract alongside the mechanisms reported in prior work.

\begin{table*}[!t]
\centering
\caption{Closest plant-facing LLM systems and the proposed framework.
Entries summarize the mechanisms and evaluation settings reported by each source.}
\label{tab:nearest}
\footnotesize
\setlength{\tabcolsep}{3pt}
\begin{tabularx}{\textwidth}{@{}p{2.7cm}XXX@{}}
\toprule
System & Strategy and write interface & Authority/context emphasized &
Reported evidence \\
\midrule
MADE \cite{huang2026made} & text strategy to structured OPC UA sequence &
operator review and authorization before execution & 40 control tasks;
virtual and physical demonstrations \\
K{\"o}hle et al.\ \cite{koehle2026llmdriven} & module agents, MCP skills,
OPC UA/MQTT & module/state-dependent tool availability & state cross-verification
for silent faults; resolution rate in a simulated six-module plant \\
Waseem et al.\ \cite{waseem2026pretrained} & LLM robot-assignment controller &
prompt constraints in a simulated line & scheduling throughput \\
\midrule
The platform & replaceable proposer above governed node/recipe execution &
channel profiles, capabilities, product/recipe limits, run context &
graded receipts, observations, attributed verdicts and batch records \\
\bottomrule
\end{tabularx}
\end{table*}

\subsection{Execution Contracts and Automation Practice}
MCP standardizes tool discovery and calls with request-scoped authorization
\cite{mcp2024}; ISA-88 and ISA-95 supply recipe and operations concepts
\cite{isa88_2010,isa95_2025}. The platform instantiates
a reduced production model and adapts run-to-run control discipline
\cite{moyne2000r2r} to agent-originated interventions. Software-agent
constraints such as AgentSpec \cite{wang2025agentspec} and CaMeL's
capability-based information-flow control \cite{debenedetti2025camel}
address tool behavior and data flow. Here those concerns meet semantic
industrial quantities, active production context, and protocol-specific
confirmation.

Electronic-record governance provides a complementary accountability
perspective: FDA guidance discusses trustworthy records and audit trails
\cite{part11}, while IEC~62443-3-3 specifies industrial system-security
requirements \cite{iec62443-3-3}. Command confirmation also has historical
telecontrol precedents, including DNP3 select-before-operate
\cite{ieee1815dnp3}. The framework brings actor identity, production context,
source-dependent admission, and retained confirmation evidence into one
supervisory workflow. Control-oriented assurance addresses state evolution
through shielding, control-barrier functions, and Simplex architectures
\cite{alshiekh2018shielding,ames2017cbf,sha2001simplex}. Those approaches
complement execution contracts that resolve which industrial operation is
authorized and what evidence its execution produces. The remaining
integration problem is to preserve those distinctions across replaceable
strategies, heterogeneous protocols, and versioned production context.
The architecture below assigns this contract to shared platform services
and retains route-specific evidence for each intervention.
